% 文献综述 + 研究动机（中文版 LaTeX 片段）
% 引用方式假定使用 natbib：\usepackage{natbib}，文中以 \citep/\citet 为主

\section{文献综述与研究动机}
\label{sec:lit_motivation_cn}

\subsection{负向经济冲击与暴力/犯罪：因果识别的经典范式}
\label{subsec:econ_shocks_crime_cn}

解释“经济下行为何诱发暴力与犯罪”的最简经济学逻辑，来自机会成本框架：当合法就业机会与预期收入下降时，个体从事非法活动的相对收益上升，从而使犯罪供给增加（\citealp{becker1968crime}）。然而，经验研究面临一个根本困难：经济状况与暴力事件往往相互影响、共同决定——冲突和犯罪既可能由衰退触发，也会反过来破坏生产与投资，使得简单相关性难以具有因果含义。

在这一点上，\citet{miguel2004economic}构成了后续研究最常被引用的“因果识别模板”。该文以降雨冲击作为收入增长的工具变量，在非洲国家面板中估计经济波动对内战爆发概率的因果效应，并得到幅度显著且稳健的负向增长冲击效应。更重要的是，作者围绕“为什么”展开机制讨论：负向冲击会扩大失业与低就业青年群体，从而降低参与暴力（内战）组织的机会成本；同时国家汲取财政与维护秩序的能力也可能因贫困而削弱。虽然其研究对象是内战而非城市街头犯罪，但它在方法论上明确了外生冲击 $\rightarrow$ 经济困境 $\rightarrow$ 暴力的识别链条，并对后续将冲击具体化、将空间尺度下沉、将结果变量细化为犯罪类型提供了可迁移的范式。

\subsection{贸易冲击与犯罪：从宏观“地方经济”到微观“社会成本”的对话}
\label{subsec:trade_crime_cn}

围绕经济周期与犯罪的关系，传统研究长期争论不休，关键原因在于：一方面犯罪测度口径差异大，另一方面经济周期指标难以被视为外生。近十余年最显著的进展是将识别重心从“宏观周期”转向“可被视为外生的局部劳动需求冲击”，其中由贸易政策/贸易自由化带来的产业冲击尤为重要。此类研究的共同点是：利用地区初始产业结构（或出口结构）作为“份额（share）”，将部门层面的外生冲击（关税削减、外需变化、进口竞争等“移动（shift）”）映射为地区层面的暴露度，从而识别贸易冲击对地方经济与社会结果的影响。

在“贸易冲击 $\rightarrow$ 犯罪”这一具体议题上，\citet{dixcarneirosoaresulyssea2018economic}（其早期工作论文版本为\citealp{dixcarneirosoaresulyssea2017economic}）是最接近你研究问题、也最适合作为综述第一段“核心靶文献”来详述的研究之一。该文利用巴西 1990--1995 年大规模单边关税削减作为准自然实验，按\citeauthor{topalova2010}--\citeauthor{kovak2013}一类做法构建地区关税冲击，并用事件研究刻画犯罪（以凶杀为主）对贸易自由化冲击的动态响应：受冲击更大的地区在改革完成后出现显著的犯罪上升，且在中期持续，长期逐步消退；同时其安慰剂与改革前趋势检验支持不存在系统性差异预趋势。更关键的是，该文没有止步于“有无效应”，而是系统比较了潜在机制变量（劳动市场、公共品供给、收入不平等/心理健康等）的动态路径，并提出利用“不同机制对冲击的响应时序差异”来对劳动市场机制的重要性进行界定（bounds）。这为后续研究提供了一个很强的写作与论证范式：\emph{先用外生冲击得到清晰的因果效应，再用机制变量的动态与异质性去识别主要传导渠道，并回应政策含义。}

与你的研究最直接对话的、发生在中国语境下的证据来自\citet{ma2025export}（你文件夹中“export slowdown and crime”）。该文使用数百万份裁判文书的文本挖掘构造城市—年份犯罪指标，并以 Bartik/shift-share 设计识别出口放缓对犯罪的影响：城市暴露度以 2010 年出口产品结构为权重，将境外需求收缩（以与中国出口结构相似国家的进口变化为代理）聚合为外生冲击。研究发现：出口放缓显著推高犯罪，并在制造业比重更高、青年/流动人口占比更高的地区更为明显。机制上，负向出口冲击降低就业与流动人口工资、恶化心理健康，而“维稳支出”等替代渠道贡献较小。该文的意义在于：它不仅给出了中国情境下“外需/贸易冲击—犯罪”的严谨证据，而且在犯罪测度上提供了可复制的“司法文本—结构化犯罪指标”路径，为在数据受限环境中开展微观犯罪经济学研究奠定了方法基础。

\subsection{中美贸易战与关税暴露：冲击构造、传导与识别策略的文献基础}
\label{subsec:tariff_exposure_cn}

你的研究以 2018 年美国对华加征关税作为外生冲击，其优势在于冲击发生时点清晰、政策来源外生性强、且与中国出口部门直接相关。因此，综述中需要同时回应两类文献：一是贸易战本身的关税归宿与价格传导证据，二是“如何把部门冲击转化为地区/网格暴露度”的方法学基础。

关于贸易战的价格传导与福利后果，\citet{amiti2019impact}在综述式写法中强调：2018 年美国加征关税主要通过提高含税进口价格影响国内消费者与进口企业，而非主要由外国出口商降价承担。更系统的量化评估来自\citet{fajgelbaum2020return}：该文估计了关税实施后的进出口数量变化与价格响应，发现进口侧与报复性关税侧均呈现近乎完全的关税传递（complete pass-through），并据此进行一般均衡下的实际收入损失核算，同时强调冲击在地区之间的分布效应（因产业集中而导致受益/受损空间分布不均）。这些证据一方面为“贸易战确实构成外生价格/需求冲击”提供经验背景，另一方面也支持在中国端把 2018 关税视作对出口与就业的外生负向冲击，从而为“关税冲击—劳动市场—犯罪”链条提供更完整的国际贸易学支撑。

关于暴露度构造与识别有效性，最典型的参照是“中国产品进口竞争冲击”一系列研究。以\citet{autor2013china}为代表的设计把地方劳动力市场视作子经济体，用初始产业结构刻画对共同冲击的差异暴露，并以对其他高收入国家的中国出口增长作为工具变量来缓解内生性担忧。尽管该类研究的结果变量主要是就业、工资与社会保障领取，而非犯罪，但它在实证策略上奠定了你所采用的“份额—移动（shift-share）”暴露度与外生冲击工具变量思路，并提供了一个重要的写作要点：\emph{暴露度来自初始结构，而“移动”来自外部冲击；识别依赖于在控制适当固定效应与趋势后，初始结构不应与未观测的差异趋势系统相关。}

近年方法学研究进一步指出，shift-share 与 Bartik 设计需要在解释与推断上更为审慎。首先，\citet{goldsmithpinkham2020bartik}强调 Bartik 工具变量的识别往往主要来自“份额（shares）”的横截面差异，其逻辑与差分中的“差异暴露”高度一致：关键不是份额与水平是否相关，而是份额是否预测了在冲击发生前后的\emph{差异趋势}。其次，\citet{adao2019shiftshare}指出传统聚类标准误可能严重低估方差，原因在于具有相似份额结构的地区残差可能相关，从而导致过度拒绝原假设；作者提出在“移动（shifters）在部门层面独立”条件下对任意跨地区残差相关均有效的推断方法。对你的研究而言，这意味着：使用连续处理强度 DID（或事件研究）时，需要更强调预趋势检验、安慰剂检验、以及与 shift-share 文献一致的稳健推断策略与聚类选择说明，从而提升识别可信度。

\subsection{城中村作为“落脚城市”：非正规住房、移民聚居与冲击脆弱性}
\label{subsec:arrival_city_cn}

上述贸易—犯罪文献的空间单位多为城市、县域或通勤区，隐含假设是“地方经济是相对均质的”。但在快速城市化语境下，外生冲击往往通过城市内部高度异质的居住—就业微环境传导。你的研究的关键创新之一，正是将“地方经济冲击”的后果下沉到城中村这一微观空间单元，并以“落脚城市”理论将这种下沉与移民转型机制连接起来。

道格·桑德斯在\citet{saunders2010arrival}中提出“落脚城市（arrival city）”概念，基于跨洲民族志材料强调：城中村/非正规住区既可能是移民融入城市、完成阶层跃迁的“跳板”，也可能在住房、联通、创业空间与公共服务不足时变为绝望与犯罪滋生的温床。该框架的价值在于，它把移民的经济机会、社会网络与空间形态捆绑在一起，提出一组可检验的预测：当外部经济冲击来临时，落脚城市的“健康程度”（低成本可达性、就业/创业承载、社会支持与教育可达性）决定其是否能够吸收冲击、抑或将冲击转化为犯罪与社会失序。

在中国语境下，城中村作为低租金、强网络的移民聚居空间，早已被城市研究与规划文献反复讨论。\citet{wang2009urbanization}以深圳为例，详细说明了城中村如何在制度与空间双重张力下形成：一方面，传统农村村落在城市扩张中被物理吞没但保留集体土地与非正规住房建设；另一方面，受户籍制度与高房价排斥，外来务工群体大量进入城中村获取可负担住房与靠近就业地的居住位置。该文强调城中村在提供廉价住房与低门槛就业方面的功能性，并提醒“大拆大建式”更新可能忽视移民福利并扰动非正规劳动力市场。对你的研究而言，这类文献提供了城中村为何会在经济冲击下更脆弱（高集中度的低技能岗位、弱保障、居住拥挤）以及为何又可能具备缓冲能力（网络、近就业、微型创业）的双重逻辑。

更进一步，住房与劳动力市场分割研究表明，移民的居住聚居与就业结果存在系统联系。\citet{liu2016segmentation}利用上海精细空间普查数据发现：农村移民居住隔离程度更高，且在就业结构上更集中于低技能、低收入岗位；同时“移民聚居区”可能通过社会网络带来更好的就业匹配与信息流动。该类发现与“落脚城市”框架相呼应：网络既可能在冲击中提供非正式互助与就业信息，发挥缓冲作用；也可能因为居住者高度同质且缺乏正式保障而导致冲击放大。由此，城中村内部与城中村之间的“健康差异”，在理论上恰恰是解释冲击异质性与缓解机制的关键。

\subsection{数据与测度前沿：司法文本、LLM 与网格级犯罪空间面板}
\label{subsec:llm_crime_cn}

“把理论下沉到城中村尺度”的最大障碍之一，是缺乏与城中村空间边界匹配的高频、细尺度犯罪数据。传统警情数据不公开且口径不一；统计年鉴等宏观指标难以反映街区层面的犯罪热点。近年“文本作为数据”的进展为中国犯罪测度打开了新路径。\citet{zhang2025llm}展示了使用大语言模型从裁判文书中抽取案件时间与地点、并结合地理编码构建全国尺度街头/邻里犯罪数据集的可行性。该研究的意义不仅在于提供数据产品，更在于示范了“非结构化司法文本 $\rightarrow$ 可用于因果推断的时空事件数据”的方法链条。

在政策评估方向，\citet{yi2023povertycrime}利用裁判文书构造县级犯罪指标，并用 DID 评估精准扶贫对刑事犯罪增长的抑制效应，机制证据指向收入增长与就业改善。这一类研究说明：司法文本可以支持经济政策的社会外部性评估，但也提示研究者必须对数据覆盖、时间分布、口径一致性与稳健性检验进行充分论证。结合你所采用的 LLM 抽取与地理编码流程，你的研究能够进一步把犯罪测度推进到“网格级别”，并与卫星影像深度学习得到的城中村分布栅格数据融合，从而在微观建成环境中实现“冲击—结果—机制”的空间对齐。

\subsection{研究缺口、研究动机与本文贡献}
\label{subsec:gap_contrib_cn}

综上，现有文献在三个方面留下了与你研究直接相关、但尚未被同时解决的缺口。
\textbf{第一}，贸易冲击与犯罪的因果链条已获得越来越多支持（如\citealp{dixcarneirosoaresulyssea2018economic,ma2025export}），但空间尺度多停留在城市/县域/通勤区，城市内部的微观异质性被均值化处理：同一座城市中哪些区域承担了更多“社会成本”，缺乏可识别的证据。%
\textbf{第二}，城中村/落脚城市文献对移民融入、非正规住房与社会网络提供了深刻机制叙事（\citealp{saunders2010arrival,wang2009urbanization}），但多为定性与描述性研究，缺乏外生冲击与严格因果推断来检验其关键命题：健康的落脚城市是否真的能在宏观冲击中“吸收”风险、降低犯罪？%
\textbf{第三}，中国语境下的犯罪数据长期受限，即便有司法文本路径，也往往难以与城市内部建成环境特征在空间上精确匹配，从而使机制检验停留在较粗粒度层面。

基于上述缺口，你的研究动机可以被组织为一个自然递进的转折：\emph{既然贸易冲击会通过劳动市场恶化推高犯罪，而城中村恰是低技能移民与非正规就业的高度集聚空间，那么，在一次外生且“干净”的关税冲击下，城中村是否会成为城市内部犯罪上升的主要承载者？更进一步，哪些“落脚城市的健康要素”能够缓解这种冲击？}

因此，本文以 2018 年美国对华加征关税作为外生冲击来源，在识别上沿用并扩展 shift-share 暴露度思路，并在推断上强调与\citet{goldsmithpinkham2020bartik}、\citet{adao2019shiftshare}一致的预趋势与稳健推断实践。与既有研究相比，你的工作至少提供三项互补贡献：%
（1）\textbf{理论贡献}：将\citeauthor{saunders2010arrival}提出的“落脚城市/城中村双重属性”从民族志叙事转化为可被量化检验的经济学命题，并把外生宏观冲击引入城中村研究框架，从而实现跨学科对话；%
（2）\textbf{数据贡献}：构建“LLM 抽取 + 地理编码”的网格级犯罪面板，并与卫星影像深度学习得到的城中村栅格分布数据融合，使得在微观建成环境尺度上进行连续处理强度 DID 成为可能；%
（3）\textbf{实证贡献}：在珠三角这一高度外向型、移民高度集聚的城市群中识别关税冲击对城中村犯罪的影响，并系统检验教育可达性、住房可负担性及微型创业空间等要素的缓解作用，从而为“如何打造更具韧性的落脚城市”提供政策相关的微观证据。

